\documentclass[a4paper, landscape]{article}

\input{include/cardLibs.tex}
\input{include/cardColours.tex}
\input{include/cardStyles.tex}

\begin{document}
\pagestyle{empty}

\begin{center}
\redcard{Tímido}{Shhh!}{Você tem a condição de tímido: não pode revelar nenhuma parte da sua carta de nenhuma forma.\\[3mm] As condições de tolo e tímido se anulam.}
\bluecard{Tímido}{Shhh!}{Você tem a condição de tímido: não pode revelar nenhuma parte da sua carta de nenhuma forma.\\[3mm] As condições de tolo e tímido se anulam.}
\redcard{Espião}{Agente duplo!}{Sua carta diz VERMELHO, mas você pertence à EQUIPE AZUL.}
\bluecard{Espião}{Agente duplo!}{Sua carta diz AZUL, mas você pertence à EQUIPE VERMELHA.}
\redcard{Reservado}{Um pouco de timidez...}{Você tem a condição de reservado: só pode compartilhar a cor. Não pode revelar a cor em público, revelar a carta em público nem compartilhar a carta.\\[3mm] As condições de tolo e reservado se anulam.}
\bluecard{Reservado}{Um pouco de timidez...}{Você tem a condição de reservado: só pode compartilhar a cor. Não pode revelar a cor em público, revelar a carta em público nem compartilhar a carta.\\[3mm] As condições de tolo e reservado se anulam.}
\redcard{Negociador}{Tudo ou nada!}{Você tem a condição de perspicaz: só pode compartilhar a carta. Não pode compartilhar a cor, revelar a cor em público nem revelar a carta em público.}
\bluecard{Negociador}{Tudo ou nada!}{Você tem a condição de perspicaz: só pode compartilhar a carta. Não pode compartilhar a cor, revelar a cor em público nem revelar a carta em público.}
\end{center}

\backpanel
\newpage

\begin{center}
\redcard{Diabo}{Mentira tem perna curta!}{Você tem a condição de mentiroso: deve sempre mentir ao falar. Isso não se aplica à comunicação não verbal.\\[3mm] As condições de honesto e mentiroso se anulam.}
\bluecard{Diabo}{Mentira tem perna curta!}{Você tem a condição de mentiroso: deve sempre mentir ao falar. Isso não se aplica à comunicação não verbal.\\[3mm] As condições de honesto e mentiroso se anulam.} 
\redcard{Anjo}{A verdade, só a verdade!}{Você tem a condição de honesto: deve sempre dizer a verdade ao falar. Isso não se aplica à comunicação não verbal.\\[3mm] As condições de mentiroso e honesto se anulam.}
\bluecard{Anjo}{A verdade, só a verdade!}{Você tem a condição de honesto: deve sempre dizer a verdade ao falar. Isso não se aplica à comunicação não verbal.\\[3mm] As condições de mentiroso e honesto se anulam.} 
\redcard{Cego}{Não espie!}{Você tem a condição de cego: faça o possível para manter os olhos fechados.\\[3mm] Não se preocupe: o jogo é curto.}
\bluecard{Cego}{Não espie!}{Você tem a condição de cego: faça o possível para manter os olhos fechados.\\[3mm] Não se preocupe: o jogo é curto.} 
\redcard{Mímico}{...}{Faça o possível para não emitir nenhum som.}
\bluecard{Mímico}{...}{Faça o possível para não emitir nenhum som.} 
\end{center}

\backpanel
\newpage

\begin{center}
\redcard{Comerciante}{Tenho uma oferta!}{Você pode trocar sua carta pela carta enterrada, sem olhar antes. Você assume todos os poderes e a equipe da nova carta.\\[3mm] Se receber outra carta de Comerciante, só poderá usar o novo poder na rodada seguinte.} 
\bluecard{Comerciante}{Tenho uma oferta!}{Você pode trocar sua carta pela carta enterrada, sem olhar antes. Você assume todos os poderes e a equipe da nova carta.\\[3mm] Se receber outra carta de Comerciante, só poderá usar o novo poder na rodada seguinte.}
\redcard{Observador}{Dê uma espiada!}{A qualquer momento, você pode olhar a carta enterrada. Oh là là!} 
\bluecard{Observador}{Dê uma espiada!}{A qualquer momento, você pode olhar a carta enterrada. Oh là là!}
\redcard{Interrogador}{Jure para mim!}{Uma vez por rodada, você pode revelar sua carta em particular a alguém da sua sala e fazer uma pergunta de sim ou não. A pessoa deve responder com sinceridade. Se realmente não souber, deve dizer isso.} 
\bluecard{Interrogador}{Jure para mim!}{Uma vez por rodada, você pode revelar sua carta em particular a alguém da sua sala e fazer uma pergunta de sim ou não. A pessoa deve responder com sinceridade. Se realmente não souber, deve dizer isso.} 
\redcard{Embaixador}{Imunidade diplomática!}{Revele sua carta imediatamente, em público, e mantenha-a revelada.\\[3mm] Você pode circular livremente entre as salas e não pode ser escolhido como refém. Porém, não pode participar de votações.\\[3mm] Você é imune: nenhuma condição pode afetá-lo.} 
\bluecard{Embaixador}{Imunidade diplomática!}{Revele sua carta imediatamente, em público, e mantenha-a revelada.\\[3mm] Você pode circular livremente entre as salas e não pode ser escolhido como refém. Porém, não pode participar de votações.\\[3mm] Você é imune: nenhuma condição pode afetá-lo.} 
\end{center}

\backpanel
\newpage

\begin{center}
\redcard{Agente}{Ninguém fica no seu caminho!}{Uma vez por rodada, revele sua carta em particular a alguém e obrigue essa pessoa a compartilhar a carta com você. Diga: ``Sou o AGENTE. Você DEVE me mostrar sua carta.''\\[3mm] Esse poder funciona até em quem normalmente não pode compartilhar a carta, como o Tímido.}
\bluecard{Agente}{Ninguém fica no seu caminho!}{Uma vez por rodada, revele sua carta em particular a alguém e obrigue essa pessoa a compartilhar a carta com você. Diga: ``Sou o AGENTE. Você DEVE me mostrar sua carta.''\\[3mm] Esse poder funciona até em quem normalmente não pode compartilhar a carta, como o Tímido.}
\redcard{Segurança}{Pode sair!}{Se sua sala tiver mais jogadores que a outra, você pode revelar sua carta em particular a alguém e dizer: ``Saia!'' Essa pessoa deve mudar de sala imediatamente.\\[3mm] Você pode fazer isso a qualquer momento, EXCETO NA ÚLTIMA RODADA.} 
\bluecard{Segurança}{Pode sair!}{Se sua sala tiver mais jogadores que a outra, você pode revelar sua carta em particular a alguém e dizer: ``Saia!'' Essa pessoa deve mudar de sala imediatamente.\\[3mm] Você pode fazer isso a qualquer momento, EXCETO NA ÚLTIMA RODADA.} 
\redcard{Vigarista}{Pode confiar em mim.}{Quando alguém aceitar compartilhar a cor com você, revele sua carta em particular em vez disso. A pessoa também deve revelar a carta dela em particular a você.}
\bluecard{Vigarista}{Pode confiar em mim.}{Quando alguém aceitar compartilhar a cor com você, revele sua carta em particular em vez disso. A pessoa também deve revelar a carta dela em particular a você.}
\redcard{Punho Vermelho}{Desvie e escape!}{A Equipe Vermelha vence se você não estiver na mesma sala que o Drone Azul ao fim do jogo.} 
\bluecard{Drone Azul}{Bip, bop, bup!}{A Equipe Azul vence se você estiver na mesma sala que o Punho Vermelho ao fim do jogo.}
\end{center}

\backpanel


\end{document}